\chapter{Column Space and Nullspace}

In the last chapter we met vector spaces and their subspaces in the abstract:
a subspace is a set of vectors that is closed under addition and scalar
multiplication, a flat sheet through the origin that you can never fall off of.
That definition is clean, but it can feel like it is floating in the air. Where
do subspaces actually \emph{come from}? Where will you meet them in practice?

The answer, and it is one of the central facts of the whole subject, is
that \textbf{every matrix carries two subspaces around with it automatically}.
A matrix $\mat{A}$ is a machine that takes input vectors $\vec{x}$ and produces
output vectors $\mat{A}\vec{x}$. Ask the two most natural questions you can ask
about any machine and you get two subspaces for free:

\begin{itemize}
  \item \textbf{What can come out?} The set of all reachable outputs is the
  \emph{column space} $\col(\mat{A})$. It lives in the output space.
  \item \textbf{What gets crushed to zero?} The set of all inputs the machine
  sends to $\vec{0}$ is the \emph{nullspace} $\nullsp(\mat{A})$. It lives in the
  input space.
\end{itemize}

These two spaces are the heart of this chapter, and arguably the heart
of the first half of the book. The column space tells you which systems
$\mat{A}\vec{x} = \vec{b}$ have a solution at all. The nullspace tells you, once
you have one solution, how many others there are. Together they are the
bookkeeping of solvability. And both of them, it turns out, fall directly out of
elimination, the very algorithm you already know.

We build the pictures first.

\section{The column space: what can come out}

\subsection{Reading a matrix--vector product as a combination of columns}

Everything in this chapter rests on one way of looking at $\mat{A}\vec{x}$ that
we have used before and will now lean on hard. When a matrix multiplies a
vector, \textbf{the output is a linear combination of the columns of
$\mat{A}$}, and the entries of $\vec{x}$ are the coefficients.

Write $\mat{A}$ with its columns named $\vec{a}_1, \vec{a}_2, \dots, \vec{a}_n$:
\[
  \mat{A} = \begin{bmatrix} \vert & \vert & & \vert \\
  \vec{a}_1 & \vec{a}_2 & \cdots & \vec{a}_n \\
  \vert & \vert & & \vert \end{bmatrix},
  \qquad
  \vec{x} = \colvec{x_1 \\ x_2 \\ \vdots \\ x_n}.
\]
Then
\[
  \mat{A}\vec{x} = x_1 \vec{a}_1 + x_2 \vec{a}_2 + \cdots + x_n \vec{a}_n.
\]
You dial in the amount $x_1$ of the first column, $x_2$ of the second, and so on,
and add up the arrows. Here it is with numbers.

\begin{example}\label{ch12:ax-as-combo}
Take
\[
  \mat{A} = \begin{bmatrix} 1 & 2 \\ 3 & 0 \\ 1 & 1 \end{bmatrix},
  \qquad \vec{x} = \colvec{4 \\ -1}.
\]
The mechanical row-by-row recipe gives
\[
  \mat{A}\vec{x} = \colvec{1\cdot 4 + 2\cdot(-1) \\ 3\cdot 4 + 0\cdot(-1)
  \\ 1\cdot 4 + 1\cdot(-1)} = \colvec{2 \\ 12 \\ 3}.
\]
But the column picture says the same thing in a more revealing way:
\[
  \mat{A}\vec{x} = 4\colvec{1 \\ 3 \\ 1} \;+\; (-1)\colvec{2 \\ 0 \\ 1}
  = \colvec{4 \\ 12 \\ 4} + \colvec{-2 \\ 0 \\ -1} = \colvec{2 \\ 12 \\ 3}.
\]
Same answer, but now you can see that the output is built out of the two
columns. Vary $\vec{x}$ over all of $\R^2$ and the output ranges over every
combination of those two columns.
\end{example}

\begin{focusbox}
$\mat{A}\vec{x}$ is a recipe for mixing the columns of $\mat{A}$. The vector
$\vec{x}$ supplies the amounts; the columns supply the directions. \emph{Hold
onto this single sentence} and the rest of the chapter unfolds from it.
\end{focusbox}

\subsection{Definition and the picture}

If $\mat{A}\vec{x}$ is always a combination of the columns, then the set of
\emph{all} outputs is exactly the set of all combinations of the columns, that
is, their span.

\begin{definition}[Column space]\label{ch12:colspace-def}
Let $\mat{A}$ be an $m\times n$ matrix with columns
$\vec{a}_1,\dots,\vec{a}_n \in \R^m$. The \textbf{column space} of $\mat{A}$,
written $\col(\mat{A})$ or $C(\mat{A})$, is the set of all linear combinations
of its columns:
\[
  \col(\mat{A}) = \spn\{\vec{a}_1, \dots, \vec{a}_n\}
  = \{\,\mat{A}\vec{x} : \vec{x} \in \R^n\,\} \subseteq \R^m.
\]
Equivalently, $\col(\mat{A})$ is the set of all vectors $\vec{b}$ for which the
system $\mat{A}\vec{x} = \vec{b}$ has at least one solution.
\end{definition}

That last sentence is worth pausing on, because it converts a geometric object
into a solvability test. Saying ``$\vec{b}$ is in the column space'' and saying
``$\mat{A}\vec{x} = \vec{b}$ is solvable'' are two phrasings of the identical
fact. The column space is precisely the set of right-hand sides you can hit.

\begin{intuitionbox}
Picture $\mat{A}$ as a machine with a panel of dials, one per column. Each dial
direction is one column vector. The column space is \emph{the set of all places
the machine can reach by turning its dials}. If your target $\vec{b}$ is inside
that reachable region, the system is solvable; if it sits outside, no setting of
the dials will ever get you there.
\end{intuitionbox}

For a $3\times 2$ matrix like the one in Example~\ref{ch12:ax-as-combo}, the two
columns are arrows in $\R^3$. As long as they do not point along the same line,
all of their combinations fill out a \emph{plane through the origin} in $\R^3$.
The output space is three-dimensional, but the reachable part is only a
two-dimensional sheet inside it. Figure~\ref{ch12:fig-colplane} shows this.

\begin{figure}[H]
\centering
\begin{tikzpicture}[scale=1.0,line join=round]
  % axes (3D-ish)
  \draw[->,spacegray] (0,0) -- (3.6,0) node[right]{$x$};
  \draw[->,spacegray] (0,0) -- (0,3.4) node[above]{$z$};
  \draw[->,spacegray] (0,0) -- (-2.1,-1.5) node[below left]{$y$};
  % the plane (a parallelogram patch of the span)
  \fill[sourceblue!18] (-2.7,-0.6) -- (1.3,1.7) -- (3.0,0.5) -- (-1.0,-1.8) -- cycle;
  \draw[sourceblue!50] (-2.7,-0.6) -- (1.3,1.7);
  \draw[sourceblue!50] (3.0,0.5) -- (-1.0,-1.8);
  \node[sourceblue!60!black] at (2.5,1.5) {$\col(\mat{A})$};
  % the two column arrows
  \draw[->,very thick,vecblue] (0,0) -- (1.3,1.7) node[above right]{$\vec{a}_1$};
  \draw[->,very thick,vecred] (0,0) -- (1.0,-0.65) node[below right]{$\vec{a}_2$};
  \fill (0,0) circle (1.6pt);
\end{tikzpicture}
\caption{Two columns of a $3\times2$ matrix are arrows in $\R^3$. All their
combinations fill a plane through the origin: the column space. Everything off
this sheet is unreachable.}
\label{ch12:fig-colplane}
\end{figure}

\subsection{Why the column space is a subspace}

We claimed $\col(\mat{A})$ is a subspace; we should not just assert it. The clean
way is to use the description $\col(\mat{A}) = \{\mat{A}\vec{x} : \vec{x}\in\R^n\}$
together with the linearity of the map $\vec{x}\mapsto\mat{A}\vec{x}$.

\begin{proposition}\label{ch12:colspace-subspace}
For any $m\times n$ matrix $\mat{A}$, the column space $\col(\mat{A})$ is a
subspace of $\R^m$.
\end{proposition}

\begin{proof}
We check the three subspace conditions for $S = \{\mat{A}\vec{x}:\vec{x}\in\R^n\}$.

\emph{Contains zero.} Take $\vec{x}=\vec{0}$. Then $\mat{A}\vec{0}=\vec{0}$, so
$\vec{0}\in S$.

\emph{Closed under addition.} Suppose $\vec{u},\vec{w}\in S$. Then $\vec{u}
= \mat{A}\vec{x}$ and $\vec{w}=\mat{A}\vec{y}$ for some inputs $\vec{x},\vec{y}$.
By linearity,
\[
  \vec{u}+\vec{w} = \mat{A}\vec{x}+\mat{A}\vec{y} = \mat{A}(\vec{x}+\vec{y}),
\]
which is again of the form $\mat{A}(\text{something})$, so $\vec{u}+\vec{w}\in S$.

\emph{Closed under scalar multiplication.} If $\vec{u}=\mat{A}\vec{x}\in S$ and
$c\in\R$, then $c\vec{u} = c(\mat{A}\vec{x}) = \mat{A}(c\vec{x})\in S$.

All three hold, so $S=\col(\mat{A})$ is a subspace. \qedhere
\end{proof}

Notice how little work that was, and notice \emph{why}. The two closure
properties are nothing but the two halves of linearity, $\mat{A}(\vec{x}+\vec{y})
= \mat{A}\vec{x}+\mat{A}\vec{y}$ and $\mat{A}(c\vec{x}) = c\,\mat{A}\vec{x}$, read
backwards. A span is always a subspace, and the column space is a span; the
proposition is really just that observation in matrix clothing.

\begin{aestheticbox}
There is something quietly beautiful here. We did not \emph{build} the column
space to be a subspace. We asked the most naive question about a matrix --
``what can come out?'' -- and the answer was forced to be a subspace by the
single fact that the matrix acts linearly. Structure is not imposed; it
\emph{emerges}. This happens again and again in linear algebra: the geometry is
already there inside the algebra, waiting for you to ask.
\end{aestheticbox}

\subsection{The shape of the column space depends on the columns}

Two columns in $\R^3$ \emph{usually} span a plane, but not always. If the second
column happens to be a multiple of the first, the two arrows lie on one line and
their span is only that line. The column space can be smaller than you expect.

\begin{example}\label{ch12:dependent-cols}
Compare
\[
  \mat{A} = \begin{bmatrix} 1 & 2 \\ 0 & 1 \\ 1 & 0 \end{bmatrix}
  \quad\text{vs.}\quad
  \mat{B} = \begin{bmatrix} 1 & 2 \\ 0 & 0 \\ 1 & 2 \end{bmatrix}.
\]
In $\mat{A}$ the columns $\colvec{1\\0\\1}$ and $\colvec{2\\1\\0}$ are not
multiples of each other, so $\col(\mat{A})$ is a plane in $\R^3$. In $\mat{B}$
the second column is exactly twice the first, $\colvec{2\\0\\2} = 2\colvec{1\\0\\1}$,
so every combination $x_1\vec{b}_1 + x_2\vec{b}_2 = (x_1+2x_2)\vec{b}_1$ is just a
scalar times $\vec{b}_1$. The column space of $\mat{B}$ is a single \emph{line}.
\end{example}

The dimension of the column space, the number of genuinely independent
directions among the columns, is so important that it has a name you have
probably already heard: the \textbf{rank}. We will study rank carefully in a
later chapter. For now keep the slogan: \emph{the column space is as big as the
columns are independent}.

\begin{trapbox}
``The column space lives in $\R^n$.'' No. Be careful which space you are in. For
an $m\times n$ matrix the columns each have $m$ entries, so $\col(\mat{A})
\subseteq \R^m$, the \emph{output} space. The nullspace, which we meet next,
lives in $\R^n$, the input space. A $3\times2$ matrix has its column space in
$\R^3$ and its nullspace in $\R^2$. Mixing these up is the single most common
beginner error in this chapter; tattoo $m$ for column space, $n$ for nullspace on
the inside of your eyelids.
\end{trapbox}

\section{The nullspace: what gets crushed to zero}

\subsection{Definition and the picture}

The column space asked about outputs. Now swing around to the input side and ask
the opposite question: which inputs does the machine destroy? Which $\vec{x}$
satisfy $\mat{A}\vec{x}=\vec{0}$?

\begin{definition}[Nullspace]\label{ch12:nullspace-def}
Let $\mat{A}$ be an $m\times n$ matrix. The \textbf{nullspace} of $\mat{A}$,
written $\nullsp(\mat{A})$ or $N(\mat{A})$, is the set of all solutions of the
homogeneous system $\mat{A}\vec{x}=\vec{0}$:
\[
  \nullsp(\mat{A}) = \{\,\vec{x}\in\R^n : \mat{A}\vec{x}=\vec{0}\,\} \subseteq \R^n.
\]
\end{definition}

The word \emph{homogeneous} just means the right-hand side is zero. The
nullspace is never empty, because $\vec{x}=\vec{0}$ is always a solution:
$\mat{A}\vec{0}=\vec{0}$. The real question is whether there are \emph{other}
solutions, nonzero inputs that the matrix flattens to the origin.

\begin{intuitionbox}
Same machine, different question. The column space was ``where can the dials take
me?'' The nullspace is ``which dial settings leave me exactly where I started, at
the origin?'' If turning the dials by $\vec{x}$ and by $\vec{x}'$ both land you at
zero, the matrix cannot tell those two settings apart at the output: it has
\emph{collapsed} a direction. The nullspace measures how much information the
matrix throws away.
\end{intuitionbox}

\subsection{Why the nullspace is a subspace}

\begin{proposition}\label{ch12:nullspace-subspace}
For any $m\times n$ matrix $\mat{A}$, the nullspace $\nullsp(\mat{A})$ is a
subspace of $\R^n$.
\end{proposition}

\begin{proof}
\emph{Contains zero.} $\mat{A}\vec{0}=\vec{0}$, so $\vec{0}\in\nullsp(\mat{A})$.

\emph{Closed under addition.} If $\mat{A}\vec{x}=\vec{0}$ and
$\mat{A}\vec{y}=\vec{0}$, then
\[
  \mat{A}(\vec{x}+\vec{y}) = \mat{A}\vec{x}+\mat{A}\vec{y} = \vec{0}+\vec{0}
  = \vec{0},
\]
so $\vec{x}+\vec{y}\in\nullsp(\mat{A})$.

\emph{Closed under scalar multiplication.} If $\mat{A}\vec{x}=\vec{0}$ and
$c\in\R$, then $\mat{A}(c\vec{x}) = c(\mat{A}\vec{x}) = c\vec{0} = \vec{0}$, so
$c\vec{x}\in\nullsp(\mat{A})$.

Hence $\nullsp(\mat{A})$ is a subspace of $\R^n$. \qedhere
\end{proof}

Geometrically: solutions of $\mat{A}\vec{x}=\vec{0}$ form a flat object through
the origin. It might be just the origin (a single point, dimension $0$), or a
line, or a plane, or more. It always passes through $\vec{0}$, and it is always
``flat'' in the subspace sense, because adding two solutions or scaling a solution
gives another solution.

\begin{trapbox}
The \emph{solution set of $\mat{A}\vec{x}=\vec{b}$ is not a subspace} unless
$\vec{b}=\vec{0}$. If $\vec{b}\neq\vec{0}$ the solution set does not contain the
origin (since $\mat{A}\vec{0}=\vec{0}\neq\vec{b}$), so it fails the very first
test. It is a flat sheet, yes, but a \emph{shifted} one that misses the origin --
what we will later call an affine set. Only the homogeneous equation gives a true
subspace. We will see in the next chapter how the shifted solution set of
$\mat{A}\vec{x}=\vec{b}$ is exactly the nullspace slid over by one particular
solution.
\end{trapbox}

\begin{figure}[H]
\centering
\begin{tikzpicture}[scale=0.98,>=Stealth]
  % LEFT: Input Space R^3
  \begin{scope}[shift={(-4.2,0)}]
    \draw[fill=gray!2, draw=gray!25, rounded corners=3pt] (-3.2,-2.7) rectangle (3.2,3.3);
    \node[draw=structurepurple!70, fill=structurepurple!12, rounded corners=2pt, font=\footnotesize\scshape, text=structurepurple!90!black] at (0,2.9) {Input Space $\R^3$};

    % 3D-ish axes
    \draw[->,spacegray] (-2.2,-0.2)--(2.4,-0.2) node[right,font=\tiny]{$x_1$};
    \draw[->,spacegray] (0,-2.0)--(0,2.4) node[above,font=\tiny]{$x_3$};
    \draw[->,spacegray] (0.8,0.6)--(-1.8,-1.4) node[below left,font=\tiny]{$x_2$};

    % Nullspace line through origin (direction squashed to 0)
    \draw[line width=1.4pt, structurepurple] (-1.8,-1.8)--(1.8,1.8);
    \node[structurepurple!90!black, font=\scriptsize] at (1.9,1.95) {$\nullsp(\mat{A})$};
    \fill[black] (0,-0.2) circle (1.8pt) node[below right,font=\tiny]{$\vec{0}$};

    % Parallel line: x_p + Null(A)
    \draw[line width=1.4pt, sourceblue] (-0.8,-2.3)--(2.8,1.3);
    \node[sourceblue!90!black, font=\scriptsize] at (2.9,1.4) {$\vec{x}_p + \nullsp(\mat{A})$};

    % Particular solution vector x_p
    \draw[->,thick,targetgreen!80!black] (0,-0.2)--(1.0,-0.5) node[midway,below,font=\tiny]{$\vec{x}_p$};
    \fill[targetgreen!80!black] (1.0,-0.5) circle (1.8pt);

    % Sample points on null line
    \fill[structurepurple] (1.0,0.8) circle (1.6pt) node[left,font=\tiny]{$\vec{n}_1$};
    \fill[structurepurple] (-1.0,-1.2) circle (1.6pt) node[left,font=\tiny]{$\vec{n}_2$};

    % Sample points on shifted line
    \fill[sourceblue] (1.0+1.0, 0.8-0.3) circle (1.6pt) node[right,font=\tiny]{$\vec{x}_p+\vec{n}_1$};
    \fill[sourceblue] (-1.0+1.0, -1.2-0.3) circle (1.6pt) node[right,font=\tiny]{$\vec{x}_p+\vec{n}_2$};

    \node[font=\tiny, spacegray, align=center] at (0,-2.3) {Entire lines map to single points};
  \end{scope}

  % Middle transition arrows
  \begin{scope}[shift={(0,0)}]
    % Null line collapse arrow
    \draw[->, thick, structurepurple] (-0.7,0.8) to[bend left=15] (0.7,0.2);
    \node[structurepurple, font=\tiny, above] at (0,0.7) {all $\mapsto \vec{0}$};

    % Shifted line collapse arrow
    \draw[->, thick, sourceblue] (-0.7,-0.6) to[bend right=15] (0.7,-1.0);
    \node[sourceblue, font=\tiny, below] at (0,-1.0) {all $\mapsto \vec{b}$};

    \node[draw=targetgreen!70, fill=targetgreen!12, rounded corners=2pt, font=\small\bfseries, text=targetgreen!85!black] at (0,2.0) {$\mat{A}\vec{x}$};
  \end{scope}

  % RIGHT: Output Space R^2
  \begin{scope}[shift={(4.2,0)}]
    \draw[fill=gray!2, draw=gray!25, rounded corners=3pt] (-3.2,-2.7) rectangle (3.2,3.3);
    \node[draw=sourceblue!70, fill=sourceblue!12, rounded corners=2pt, font=\footnotesize\scshape, text=sourceblue!90!black] at (0,2.9) {Output Space $\R^2$};

    % 2D axes
    \draw[->,spacegray] (-2.4,0)--(2.5,0) node[right,font=\tiny]{$y_1$};
    \draw[->,spacegray] (0,-2.2)--(0,2.4) node[above,font=\tiny]{$y_2$};

    % Column space: reachable plane (here all of R^2 or a line)
    \draw[very thick, targetgreen!80!black] (-2.2,-1.1)--(2.2,1.1);
    \node[targetgreen!80!black, font=\scriptsize] at (2.1,1.3) {$\col(\mat{A})$};

    % Collapsed origin
    \fill[structurepurple] (0,0) circle (2.5pt);
    \node[structurepurple!90!black, font=\scriptsize, above left] at (0,0.05) {$\vec{0} = \mat{A}(\nullsp)$};

    % Collapsed target vector b
    \fill[sourceblue] (1.2,0.6) circle (2.5pt);
    \node[sourceblue!90!black, font=\scriptsize, below right] at (1.2,0.55) {$\vec{b} = \mat{A}\vec{x}_p$};
    \draw[->,thick,sourceblue] (0,0)--(1.2,0.6);

    \node[font=\tiny, spacegray, align=center] at (0,-2.3) {Information along $\nullsp(\mat{A})$ is lost};
  \end{scope}
\end{tikzpicture}
\caption{The geometry of nullspace collapse. The matrix $\mat{A}$ squashes the entire purple line
$\nullsp(\mat{A})$ through the origin down to the single point $\vec{0}$ in the output space.
Consequently, every parallel shifted fiber $\vec{x}_p + \nullsp(\mat{A})$ (blue) collapses
to the exact same output vector $\vec{b} = \mat{A}\vec{x}_p$. This explains why the complete solution
to $\mat{A}\vec{x}=\vec{b}$ is an entire affine line: any two inputs that differ by a nullspace
vector produce identical outputs.}
\label{ch12:fig-two-spaces}
\end{figure}

\section{Finding the nullspace by elimination}

So far the nullspace is a definition. To actually \emph{find} it, to write down
all the vectors $\vec{x}$ with $\mat{A}\vec{x}=\vec{0}$, we use the one tool we
have polished to a shine: elimination. The key fact is that elimination
does not disturb the nullspace at all.

\subsection{Elimination preserves the nullspace}

\begin{lemma}\label{ch12:row-ops-preserve-null}
If $\mat{R}$ is obtained from $\mat{A}$ by elementary row operations (row swaps,
scaling a row, adding a multiple of one row to another), then $\mat{A}$ and
$\mat{R}$ have exactly the same nullspace.
\end{lemma}

\begin{proof}[Why]
Each elementary row operation can be undone by another elementary row operation,
so the rows of $\mat{R}$ are combinations of the rows of $\mat{A}$ and vice
versa. The equation $\mat{A}\vec{x}=\vec{0}$ says every row of $\mat{A}$ dotted
with $\vec{x}$ gives $0$. If that holds, then any combination of the rows also
dots with $\vec{x}$ to give $0$, so $\mat{R}\vec{x}=\vec{0}$ too; and because the
operations are reversible, the converse holds as well. The solution set never
changes. (This is special to row operations and the right-hand side $\vec
0$: column operations would \emph{not} preserve the nullspace.)
\end{proof}

This is the green light we needed. To find $\nullsp(\mat{A})$ we may freely reduce
$\mat{A}$ to its row echelon form, or all the way to the \emph{reduced} row
echelon form, and read off the answer there, where it is easy.

\subsection{Pivots, free variables, and special solutions}

Run elimination on $\mat{A}$. Some columns end up with pivots (the first nonzero
entry in their row); these are the \textbf{pivot columns} and their variables are
\textbf{pivot variables}. The columns \emph{without} a pivot correspond to
\textbf{free variables}. The name is exactly right: you are free to set those
variables to anything you like, and the pivot variables are then forced. Each
free variable gives you one ``knob,'' and turning one knob at a time produces one
\textbf{special solution}. The whole nullspace is the span of these special
solutions.

\begin{focusbox}
\textbf{The recipe.}
\begin{enumerate}
  \item Reduce $\mat{A}$ to reduced row echelon form $\mat{R}$.
  \item Identify pivot columns and free columns. The number of free columns is
  $n - (\text{number of pivots})$.
  \item For each free variable, set it to $1$ and all other free variables to
  $0$, then solve for the pivot variables. That vector is a special solution.
  \item $\nullsp(\mat{A})$ is the span of all the special solutions.
\end{enumerate}
\end{focusbox}

We grind through this on a concrete matrix until it is muscle memory.

\begin{example}[A full nullspace computation]\label{ch12:big-null}
Find the nullspace of
\[
  \mat{A} = \begin{bmatrix}
    1 & 2 & 2 & 4 \\
    1 & 2 & 3 & 6 \\
    2 & 4 & 5 & 10
  \end{bmatrix}.
\]

\textbf{Step 1: eliminate.} Subtract row 1 from row 2, and subtract $2\times$ row
1 from row 3:
\[
  \begin{bmatrix}
    1 & 2 & 2 & 4 \\
    0 & 0 & 1 & 2 \\
    0 & 0 & 1 & 2
  \end{bmatrix}.
\]
Now subtract the new row 2 from row 3:
\[
  \begin{bmatrix}
    1 & 2 & 2 & 4 \\
    0 & 0 & 1 & 2 \\
    0 & 0 & 0 & 0
  \end{bmatrix}.
\]
Finally clear above the second pivot (subtract $2\times$ row 2 from row 1) to
reach reduced row echelon form:
\[
  \mat{R} = \begin{bmatrix}
    1 & 2 & 0 & 0 \\
    0 & 0 & 1 & 2 \\
    0 & 0 & 0 & 0
  \end{bmatrix}.
\]

\textbf{Step 2: pivots and free variables.} The pivots sit in columns $1$ and
$3$. So $x_1$ and $x_3$ are pivot variables, while $x_2$ and $x_4$ are
\emph{free}. There are $4-2=2$ free variables, so we expect a two-dimensional
nullspace, a plane in $\R^4$.

\textbf{Step 3: special solutions.} The reduced rows say
\[
  x_1 + 2x_2 = 0 \quad\Longrightarrow\quad x_1 = -2x_2,
  \qquad
  x_3 + 2x_4 = 0 \quad\Longrightarrow\quad x_3 = -2x_4.
\]
First special solution $\vec{s}_1$: set $x_2 = 1$, $x_4 = 0$. Then
$x_1 = -2$, $x_3 = 0$:
\[
  \vec{s}_1 = \colvec{-2 \\ 1 \\ 0 \\ 0}.
\]
Second special solution $\vec{s}_2$: set $x_2 = 0$, $x_4 = 1$. Then $x_1 = 0$,
$x_3 = -2$:
\[
  \vec{s}_2 = \colvec{0 \\ 0 \\ -2 \\ 1}.
\]

\textbf{Step 4: the nullspace.} Every solution is a combination of these:
\[
  \nullsp(\mat{A}) = \left\{\, c_1 \colvec{-2\\1\\0\\0} + c_2\colvec{0\\0\\-2\\1}
  : c_1, c_2 \in\R \,\right\}.
\]

\textbf{Check.} Multiply $\mat{A}\vec{s}_1$:
$1(-2)+2(1)+2(0)+4(0) = 0$, $1(-2)+2(1)+3(0)+6(0)=0$, $2(-2)+4(1)+5(0)+10(0)=0$.
All zero, as required. (You should check $\vec{s}_2$ yourself; it is good for the
soul.)
\end{example}

\begin{labbox}
Find the nullspace of $\mat{A}=\begin{bmatrix}1 & 3 & 1\\2 & 6 & 4\end{bmatrix}$
by hand. Reduce, spot the pivot in column 1 and another in column 3, notice $x_2$
is free, and produce the one special solution. You should get $\nullsp(\mat{A}) =
\spn\{(-3,1,0)\}$, a line in $\R^3$. Then \emph{verify} by computing
$\mat{A}\,(-3,1,0)\T$ and confirming it is $\vec 0$. Doing this once slowly is
worth reading three pages.
\end{labbox}

\subsection{Why the special solutions span the whole nullspace}

We asserted that the special solutions capture \emph{every} nullspace vector. Let
us see why; the reason is clean, and it is the same reason every time.

After reduction, the free variables can be \emph{anything}. Take any solution
$\vec{x}$ of $\mat{R}\vec{x}=\vec{0}$ and read off its free-variable entries; call
them $f_1, f_2, \dots$ (in Example~\ref{ch12:big-null} these are $x_2$ and
$x_4$). Now form the combination
\[
  f_1\vec{s}_1 + f_2\vec{s}_2 + \cdots,
\]
where $\vec{s}_j$ is the special solution for the $j$-th free variable. By
construction $\vec{s}_j$ has a $1$ in free-slot $j$ and a $0$ in the other free
slots, so this combination reproduces \emph{exactly} the free entries
$f_1,f_2,\dots$ of $\vec{x}$. And since each $\vec{s}_j$ already satisfies the
homogeneous equation, so does the combination, which means its pivot entries are
forced to be the unique correct ones, matching $\vec{x}$. Two solutions with the
same free entries must be the same solution. Therefore
\[
  \vec{x} = f_1\vec{s}_1 + f_2\vec{s}_2 + \cdots.
\]
Every nullspace vector is a combination of the special solutions: they span the
nullspace. (They are also independent, since each owns a free slot the others
zero out, so they form a \emph{basis}, which is why $\dim\nullsp(\mat{A})$ equals
the number of free variables. More on bases soon.)

\begin{formalbox}
\textbf{Counting solutions of $\mat{A}\vec{x}=\vec{0}$.} For an $m\times n$ matrix
$\mat{A}$ with $r$ pivots after elimination:
\begin{itemize}
  \item the number of free variables is $n - r$;
  \item $\dim\nullsp(\mat{A}) = n - r$, with the $n-r$ special solutions as a
  basis;
  \item if $r = n$ (a pivot in every column, no free variables) then
  $\nullsp(\mat{A}) = \{\vec{0}\}$: the only solution is the trivial one;
  \item if $r < n$ then there are infinitely many solutions, filling an
  $(n-r)$-dimensional subspace.
\end{itemize}
In particular, if a system has more unknowns than equations ($n > m \ge r$),
there must be free variables, so $\mat{A}\vec{x}=\vec{0}$ always has nonzero
solutions. \emph{A wide homogeneous system can never have only the zero
solution.}
\end{formalbox}

\begin{example}[Nullspace is only the origin]\label{ch12:trivial-null}
\[
  \mat{A} = \begin{bmatrix} 1 & 2 \\ 0 & 3 \\ 1 & 1 \end{bmatrix}
  \;\xrightarrow{\text{eliminate}}\;
  \begin{bmatrix} 1 & 2 \\ 0 & 3 \\ 0 & 0 \end{bmatrix}.
\]
Both columns have pivots ($r=2=n$), so there are no free variables. The only
solution of $\mat{A}\vec{x}=\vec{0}$ is $\vec{x}=\vec{0}$, and
$\nullsp(\mat{A})=\{\vec 0\}$. The columns are independent; the matrix crushes
nothing.
\end{example}

\subsection{A rank-one example tying both spaces together}

\begin{example}[One matrix, both spaces]\label{ch12:nullspace-line}
\[
  \mat{A} = \begin{bmatrix} 1 & 2 \\ 2 & 4 \end{bmatrix}.
\]
The second column is twice the first, so the column space is the line
$\spn\{(1,2)\}$ in $\R^2$ -- everything the matrix produces lies on that line.
Now the nullspace. Eliminate: subtract $2\times$ row 1 from row 2 to get
$\begin{bmatrix}1 & 2\\0 & 0\end{bmatrix}$. Column 1 has the only pivot; $x_2$ is
free. The pivot row says $x_1 + 2x_2 = 0$, i.e. $x_1 = -2x_2$. The special
solution (set $x_2=1$) is $(-2,1)$, so
\[
  \nullsp(\mat{A}) = \spn\left\{\colvec{-2\\1}\right\},
\]
a line in $\R^2$. Geometrically the matrix squashes the whole plane onto a single
line (the column space), and the direction $(-2,1)$ is exactly the direction
squashed to the origin (the nullspace). Figure~\ref{ch12:fig-two-spaces}
illustrates this collapse: one input dimension dies, and correspondingly the
output is only one-dimensional.
\end{example}

The collapse is easiest to feel by watching what the matrix does to a whole
patch of the plane. Feed the unit square into $\mat{A}=\begin{bmatrix}1&2\\2&4
\end{bmatrix}$ and it does not come out as a tilted parallelogram, the way an
invertible matrix would leave it. It comes out \emph{flattened onto a single
line}: the column space. Every point of the square is pushed onto
$\spn\{(1,2)\}$, and the entire nullspace direction $(-2,1)$ is what gets
ironed flat. Figure~\ref{ch12:fig-collapse} shows the squash.

\begin{figure}[H]
\centering
\begin{tikzpicture}[scale=0.92,>=Stealth]
  % Left panel: Input plane
  \begin{scope}[shift={(-3.8,0)}]
    \draw[fill=gray!2, draw=gray!25, rounded corners=3pt] (-2.2,-1.4) rectangle (3.0,3.6);
    \node[draw=structurepurple!70, fill=structurepurple!12, rounded corners=2pt, font=\footnotesize\scshape, text=structurepurple!90!black] at (0.4,3.2) {Domain $\R^2$: Input Square};
    \draw[->,spacegray] (-0.8,0) -- (2.6,0) node[right,font=\tiny]{$x_1$};
    \draw[->,spacegray] (0,-0.8) -- (0,2.6) node[above,font=\tiny]{$x_2$};

    % unit square
    \fill[sourceblue!18] (0,0) -- (1.2,0) -- (1.2,1.2) -- (0,1.2) -- cycle;
    \draw[sourceblue!70,thick] (0,0) -- (1.2,0) -- (1.2,1.2) -- (0,1.2) -- cycle;
    \node[sourceblue!80!black, font=\scriptsize] at (0.6,0.6) {Area $= 1$};

    % nullspace direction (-2,1)
    \draw[very thick,structurepurple,dashed] (1.6,-0.8) -- (-1.6,0.8);
    \node[structurepurple!90!black, font=\scriptsize, below] at (1.6,-0.85) {$\nullsp(\mat{A})$ direction};
    \draw[->,very thick,structurepurple] (0,0) -- (-1.2,0.6)
      node[above left,font=\scriptsize]{$(-2,1)$};
    \fill (0,0) circle (1.6pt) node[below left,font=\tiny]{$\vec{0}$};
  \end{scope}

  % Middle action arrow
  \draw[->,very thick,targetgreen!80!black] (-0.5,1.2) -- (0.8,1.2)
    node[midway,above,font=\small\bfseries]{$\mat{A}$};

  % Right panel: Output plane
  \begin{scope}[shift={(3.8,0)}]
    \draw[fill=gray!2, draw=gray!25, rounded corners=3pt] (-1.8,-1.4) rectangle (3.4,3.6);
    \node[draw=sourceblue!70, fill=sourceblue!12, rounded corners=2pt, font=\footnotesize\scshape, text=sourceblue!90!black] at (0.8,3.2) {Codomain $\R^2$: Collapsed Line};
    \draw[->,spacegray] (-0.8,0) -- (2.8,0) node[right,font=\tiny]{$y_1$};
    \draw[->,spacegray] (0,-0.8) -- (0,2.8) node[above,font=\tiny]{$y_2$};

    % column space line y = 2x
    \draw[sourceblue!45,thin] (-0.5,-1.0) -- (1.5,3.0);
    \node[sourceblue!75!black, font=\scriptsize, right] at (1.5,2.9) {$\col(\mat{A}): y_2=2y_1$};

    % image of unit square: flattened segment
    \draw[line width=2.4pt,vecblue] (0,0) -- (1.3,2.6);
    \draw[->,very thick,vecblue] (0,0) -- (0.8,1.6) node[right,font=\scriptsize]{$\vec{a}_1=\colvec{1\\2}$};
    \fill[black] (0,0) circle (1.8pt) node[below right,font=\tiny]{$\vec{0}$};
    \node[font=\scriptsize, align=center, sourceblue!90!black] at (0.8,-0.9) {Area collapsed to $0$\\(flat segment)};
  \end{scope}
\end{tikzpicture}
\caption{The rank-one matrix $\mat{A}=\protect\begin{bmatrix}1&2\\2&4\protect
\end{bmatrix}$ flattens the two-dimensional unit square onto a one-dimensional segment
of the line $\col(\mat{A})=\spn\{(1,2)\}$. The entire nullspace direction $(-2,1)$ is
squashed to zero, destroying all 2D area.}
\label{ch12:fig-collapse}
\end{figure}

\begin{walkbox}
Take the matrix $\mat{A}=\begin{bmatrix}1&1\\1&1\end{bmatrix}$ out for a walk.
Without writing anything down, answer: what is its column space (a line? which
line?), and what direction does it crush to zero? Then ask the deeper one: the
nullspace direction $(1,-1)$ and the column space direction $(1,1)$ are
perpendicular. Is that a coincidence, or is something structural going on between
the rows of $\mat{A}$ and its nullspace? Hold that question; the chapter on the
four fundamental subspaces will reward it.
\end{walkbox}

\section{The two spaces as a solvability ledger}

Step back and see what we have built. The column space and the nullspace
together answer, completely, the two questions you can ask about
$\mat{A}\vec{x}=\vec{b}$:

\begin{center}
\begin{tabular}{ll}
\toprule
\textbf{Question} & \textbf{Answered by} \\
\midrule
Does a solution \emph{exist}? & Is $\vec{b}\in\col(\mat{A})$? \\
If one exists, is it \emph{unique}? & Is $\nullsp(\mat{A})=\{\vec 0\}$? \\
\bottomrule
\end{tabular}
\end{center}

Existence is the column space talking. Uniqueness is the nullspace talking. Here
is the uniqueness half spelled out, because it is one of the most useful small
facts in the subject.

\begin{theorem}[Nullspace controls uniqueness]\label{ch12:uniqueness}
Suppose $\mat{A}\vec{x}=\vec{b}$ has at least one solution $\vec{x}_p$. Then the
complete set of solutions is
\[
  \{\,\vec{x}_p + \vec{n} : \vec{n}\in\nullsp(\mat{A})\,\}.
\]
In particular the solution is unique if and only if $\nullsp(\mat{A})=\{\vec 0\}$.
\end{theorem}

\begin{proof}
First, every such vector solves the system: $\mat{A}(\vec{x}_p+\vec{n})
= \mat{A}\vec{x}_p + \mat{A}\vec{n} = \vec{b} + \vec{0} = \vec{b}$.

Conversely, let $\vec{x}$ be \emph{any} solution. Then
\[
  \mat{A}(\vec{x}-\vec{x}_p) = \mat{A}\vec{x} - \mat{A}\vec{x}_p
  = \vec{b}-\vec{b} = \vec{0},
\]
so $\vec{n} := \vec{x}-\vec{x}_p$ lies in the nullspace, and $\vec{x}
= \vec{x}_p + \vec{n}$ has the claimed form. Thus the solution set is exactly the
nullspace shifted by $\vec{x}_p$. It collapses to the single point $\vec{x}_p$
precisely when the nullspace is $\{\vec 0\}$. \qedhere
\end{proof}

This is the theorem the next chapter is built on: the complete solution of
$\mat{A}\vec{x}=\vec{b}$ is ``one particular solution plus all of the
nullspace.'' The nullspace is the freedom; the particular solution is one place
to start. Two solutions can differ only by a nullspace vector, never by anything
else. Figure~\ref{ch12:fig-affine} draws the picture: the solution set is the
nullspace line, picked up off the origin and laid down through $\vec{x}_p$.

\begin{figure}[H]
\centering
\begin{tikzpicture}[scale=1.0,line join=round]
  \draw[->,spacegray] (-2.6,0) -- (3.2,0) node[right]{\small $x_1$};
  \draw[->,spacegray] (0,-1.6) -- (0,3.0) node[above]{\small $x_2$};
  % nullspace line through origin, direction (2,1)
  \draw[very thick,structurepurple] (-2.4,-1.2) -- (2.4,1.2);
  \node[structurepurple] at (2.7,1.25) {\small $\nullsp(\mat{A})$};
  % solution set: parallel line through x_p = (0.5,2)
  \draw[very thick,sourceblue] (-1.9,0.8) -- (2.9,3.2);
  \node[sourceblue!70!black] at (2.9,3.45) {\small solution set};
  % the shift vector
  \draw[->,thick,targetgreen] (0,0) -- (0.5,2);
  \node[targetgreen!60!black,right] at (0.25,1.15) {\small $\vec{x}_p$};
  \fill[sourceblue] (0.5,2) circle (1.8pt);
  % a nullspace vector added on
  \draw[->,thick,structurepurple,dashed] (0.5,2) -- (2.1,2.8);
  \node[structurepurple] at (1.9,3.1) {\small $+\,\vec{n}$};
  \fill[sourceblue] (2.1,2.8) circle (1.5pt);
  \fill (0,0) circle (1.7pt);
  \node[below left] at (0,0) {\small $\vec 0$};
\end{tikzpicture}
\caption{Theorem~\ref{ch12:uniqueness} in one picture. The homogeneous solutions
$\nullsp(\mat{A})$ form a line through the origin (purple). Adding one particular
solution $\vec{x}_p$ slides that whole line off the origin to give the solution
set of $\mat{A}\vec{x}=\vec{b}$ (blue): a parallel, \emph{shifted} copy. Any two
solutions differ by a nullspace vector $\vec{n}$. The set is a single point
exactly when the nullspace is just $\{\vec 0\}$.}
\label{ch12:fig-affine}
\end{figure}

\begin{strangbox}
If you have not yet read Strang: when you get to his lectures, notice that he
introduces $C(\mat{A})$ and $N(\mat{A})$ almost in the same breath, and that he
insists on the column picture of $\mat{A}\vec{x}$ from day one. The reason is
exactly Theorem~\ref{ch12:uniqueness}: these two subspaces are how he organizes
\emph{all} of $\mat{A}\vec{x}=\vec{b}$. If you have already finished Strang: you
will recognize this chapter as the front half of his ``four fundamental
subspaces.'' We have deliberately met the column space and nullspace alone first,
so that when their two partners (the row space and the left nullspace) arrive, you
already have half the picture and can feel the symmetry snap into place.
\end{strangbox}

\section{A modern application and a connection}

\begin{applicationbox}
\textbf{Nullspace and column space in machine learning.} Stack your data as a
matrix $\mat{X}$ with one feature per column. Fitting a linear model means
finding weights $\vec{w}$ so that $\mat{X}\vec{w}$ approximates the target. Two
things from this chapter bite immediately.

\emph{Redundant features make a nontrivial nullspace.} If one feature is a linear
combination of others -- say ``height in inches'' alongside ``height in
centimeters,'' or a one-hot encoding that includes every category plus an
all-ones intercept -- then $\mat{X}$ has a nonzero nullspace. By
Theorem~\ref{ch12:uniqueness} the best-fit weights are then \emph{not unique}:
you can add any nullspace vector to $\vec{w}$ and the predictions
$\mat{X}\vec{w}$ do not change at all. This is the linear-algebra heart of
\emph{multicollinearity}, and it is why software either warns you, drops a
column, or adds regularization to pick one answer out of the infinite flat sheet
of equally good ones.

\emph{The column space is the set of achievable predictions.} No matter how you
choose $\vec{w}$, the prediction $\mat{X}\vec{w}$ lives in $\col(\mat{X})$. If the
true targets do not lie in that span, you simply cannot fit them exactly; the
best you can do is land at the closest point of the column space. That closest
point is the orthogonal projection onto $\col(\mat{X})$, and it is exactly what
least-squares regression computes -- the subject of a chapter to come. So the
column space is not an abstraction here; it is the literal menu of predictions
your model is allowed to make.
\end{applicationbox}

\begin{connectionbox}
\textbf{Kernel and image, everywhere in mathematics.} The two words for our two
spaces in the wider mathematical world are \emph{kernel} (the nullspace) and
\emph{image} (the column space), and they appear for \emph{any} linear map, not
just matrices.

In calculus, differentiation $D: f \mapsto f'$ is linear. Its kernel is the set
of functions with $f' = 0$ -- the constants. That is why every indefinite
integral comes with a ``$+\,C$'': the $+C$ is precisely the nullspace of $D$, and
Theorem~\ref{ch12:uniqueness} is the reason the general antiderivative is ``one
particular antiderivative plus the kernel.'' The structure you learned for
$\mat{A}\vec{x}=\vec{b}$ is \emph{the same structure} as ``solve $f' = g$.''

The same pattern governs solutions of linear differential equations (general
solution $=$ particular solution $+$ homogeneous solutions), and it is the seed
of the rank--nullity theorem and, much deeper, the first isomorphism theorem in
group and module theory. Kernel measures collapse; image measures reach. Once you
see that one pair of ideas running underneath linear systems, calculus, and
abstract algebra alike, the landscape of mathematics gets a good deal smaller and
friendlier.
\end{connectionbox}

\section{Reading the column space off the echelon form (a caution)}

Elimination gave us the nullspace cleanly. It is tempting to expect it to hand us
the column space just as cleanly, but there is a subtlety worth a warning, because
it trips up almost everyone once.

\begin{trapbox}
\textbf{The pivot columns tell you \emph{which} original columns form a basis for
$\col(\mat{A})$ -- but you must take them from $\mat{A}$, not from $\mat{R}$.}
Row operations \emph{change} the column space (they mix the entries within each
column), so $\col(\mat{R})\neq\col(\mat{A})$ in general. What \emph{is} preserved
is the \emph{pattern} of dependence among columns: column $j$ of $\mat{A}$ is a
combination of earlier columns exactly when column $j$ of $\mat{R}$ is. So you
reduce to find \emph{which} columns are pivot columns, then go back to the
\emph{original} $\mat{A}$ and take those columns as your basis.
\end{trapbox}

\begin{example}\label{ch12:colspace-basis}
Return to Example~\ref{ch12:big-null}, where
\[
  \mat{A} = \begin{bmatrix} 1 & 2 & 2 & 4 \\ 1 & 2 & 3 & 6 \\ 2 & 4 & 5 & 10
  \end{bmatrix}
  \xrightarrow{\text{rref}}
  \mat{R} = \begin{bmatrix} 1 & 2 & 0 & 0 \\ 0 & 0 & 1 & 2 \\ 0 & 0 & 0 & 0
  \end{bmatrix}.
\]
Pivots are in columns $1$ and $3$. So a basis for $\col(\mat{A})$ is columns $1$
and $3$ \emph{of $\mat{A}$}:
\[
  \col(\mat{A}) = \spn\left\{\colvec{1\\1\\2},\ \colvec{2\\3\\5}\right\}.
\]
Columns $2$ and $4$ are redundant. Indeed $\mat{R}$ shows column $2$ is $2\times$
column $1$ and column $4$ is $2\times$ column $3$, and you can check the same
multiples hold in $\mat{A}$: $\colvec{2\\2\\4}=2\colvec{1\\1\\2}$ and
$\colvec{4\\6\\10}=2\colvec{2\\3\\5}$. The column space is a plane in $\R^3$,
spanned by those two surviving columns. Notice the same number, $r=2$, counts both the pivots \emph{and} the dimension of the column space:
each pivot column isolates a genuinely new direction, providing the geometric foundation of the rank theorem.
\end{example}

Notice the bookkeeping: $r = 2$ pivots gave $n - r = 2$ free variables and
a $2$-dimensional nullspace, while the $r = 2$ pivot columns gave a
$2$-dimensional column space. Two pivots, doing double duty. The conservation law
hiding here, $\dim\col(\mat{A}) + \dim\nullsp(\mat{A}) = n$, is the
\emph{rank--nullity theorem}, and we will prove it in full when we have bases and
dimension firmly in hand. For now, note that the same elimination, run once,
reveals \emph{both} subspaces of the matrix.

\begin{historybox}
Throughout the eighteenth and early nineteenth centuries, mathematicians viewed a system of linear equations purely as a mechanical puzzle: given a set of numerical constraints, perform algebraic substitutions until each variable is isolated. There was no concept of a matrix as an active operator, nor any notion that the inputs and outputs lived in distinct geometric spaces.

In the 1870s, the French mathematician Camille Jordan and the German mathematician Ferdinand Georg Frobenius revolutionized this view. They realized that every linear substitution decomposes the universe into two complementary subspaces: the \emph{image} (what we now call the column space $\col(\mat{A})$) and the \emph{kernel} (what we call the nullspace $\nullsp(\mat{A})$). The original question Frobenius and Jordan settled was clarifying exactly when a system $\mat{A}\vec{x}=\vec{b}$ can be solved, and how much ambiguity remains when solutions exist. Before their work, algebraists often confused the equations that constrain a system with the variables that generate its solutions. Frobenius showed that solvability is a reachability test: $\vec{b}$ must live in the subspace spanned by the columns. Uniqueness, conversely, depends entirely on whether the kernel contains anything beyond the zero vector.

A persistent historical trap was the confusion between row operations and column operations. Because elementary row operations preserve all algebraic dependencies between columns, mathematicians found they could read the dimension of both spaces directly off the echelon form $\mat{R}$. But early practitioners frequently made the fatal error of assuming that the column space of $\mat{A}$ was spanned by the pivot columns of $\mat{R}$. It took careful axiomatization to cement the distinction: row operations preserve the nullspace because they do not change the solution set of $\mat{A}\vec{x}=\vec{0}$, but they violently tilt the column space by mixing row coordinates.
\end{historybox}

\section{Summary}

\begin{focusbox}
\textbf{What to carry out of this chapter.}
\begin{itemize}
  \item $\mat{A}\vec{x}$ is a linear combination of the columns of $\mat{A}$, with
  $\vec{x}$ as the coefficients. Everything follows from this.
  \item $\col(\mat{A})\subseteq\R^m$ is the span of the columns: all reachable
  outputs, equivalently all $\vec{b}$ making $\mat{A}\vec{x}=\vec{b}$ solvable.
  \item $\nullsp(\mat{A})\subseteq\R^n$ is the set of all $\vec{x}$ with
  $\mat{A}\vec{x}=\vec{0}$: the inputs crushed to zero.
  \item Both are subspaces, forced by the linearity of $\mat{A}$.
  \item Find the nullspace by elimination: reduce, find free variables, build one
  special solution per free variable, and span them. $\dim\nullsp(\mat{A})
  = n - r$.
  \item Find a basis for the column space by taking the \emph{pivot columns of the
  original $\mat{A}$}. $\dim\col(\mat{A}) = r$.
  \item Column space governs \emph{existence}; nullspace governs \emph{uniqueness}.
  The complete solution of $\mat{A}\vec{x}=\vec{b}$ is one particular solution
  plus all of the nullspace.
\end{itemize}
\end{focusbox}

\begin{connectionbox}
The column space and nullspace give us two clean geometric answers to any linear system $\mat{A}\vec{x}=\vec{b}$. Existence is completely answered by whether $\vec{b} \in \col(\mat{A})$; uniqueness is completely answered by whether $\nullsp(\mat{A}) = \{\vec{0}\}$. One pass of Gaussian elimination produces both subspaces simultaneously.

Yet knowing that $\vec{b}$ lies in $\col(\mat{A})$ guarantees only that a solution exists; it does not hand us a formula for that solution. Furthermore, what happens when $\vec{b} \notin \col(\mat{A})$? In real-world data and physical measurements, noise almost always pushes $\vec{b}$ outside the column space, rendering exact solutions impossible. In Chapter~13, we assemble these two spaces into the complete solution $\vec{x} = \vec{x}_p + \vec{x}_n$. Later, in Chapters 18 and 19, we learn how to overcome unsolvable systems by projecting $\vec{b}$ orthogonally onto $\col(\mat{A})$ to obtain the best possible least-squares fit.
\end{connectionbox}

\section{Exercises}

\begin{exercise}[Warmup: combination of columns]\label{ch12:ex-combo}
For $\mat{A}=\begin{bmatrix}2 & 0\\1 & 3\\0 & 1\end{bmatrix}$ and
$\vec{x}=\colvec{3\\-2}$, compute $\mat{A}\vec{x}$ two ways: by the row recipe,
and by writing it as $3(\text{col }1) + (-2)(\text{col }2)$. Confirm the answers
agree, and state which space the answer lives in, $\R^2$ or $\R^3$.
\end{exercise}

\begin{exercise}[Is $\vec b$ reachable?]\label{ch12:ex-reachable}
Let $\mat{A}=\begin{bmatrix}1 & 1\\1 & 2\\1 & 3\end{bmatrix}$. Determine whether
$\vec{b}=\colvec{1\\0\\2}$ lies in $\col(\mat{A})$ by trying to solve
$\mat{A}\vec{x}=\vec{b}$. If it is unreachable, give one equation that every
vector in $\col(\mat{A})$ must satisfy (hint: the three entries of any column
combination obey a linear relation).
\end{exercise}

\begin{exercise}[Nullspace by elimination]\label{ch12:ex-null}
Find the nullspace of
$\mat{A}=\begin{bmatrix}1 & 2 & 1 & 1\\2 & 4 & 3 & 4\end{bmatrix}$.
Identify the pivot and free variables, write down each special solution, and
state $\dim\nullsp(\mat{A})$. Verify one special solution by direct multiplication.
\end{exercise}

\begin{exercise}[Trivial nullspace]\label{ch12:ex-trivial}
For which value(s) of $t$ does $\mat{A}=\begin{bmatrix}1 & 2\\3 & t\end{bmatrix}$
have nullspace equal to $\{\vec 0\}$? For the value of $t$ that fails, describe
the nullspace as a line and give its direction.
\end{exercise}

\begin{exercise}[Column-space basis]\label{ch12:ex-colbasis}
Let $\mat{A}=\begin{bmatrix}1 & 2 & 3\\2 & 4 & 7\\1 & 2 & 4\end{bmatrix}$. Reduce
to find the pivot columns, then give a basis for $\col(\mat{A})$ taken from the
columns of $\mat{A}$. What is $\dim\col(\mat{A})$, and what is $\dim\nullsp(\mat
A)$? Check that they add to $3$.
\end{exercise}

\begin{exercise}[Particular plus nullspace]\label{ch12:ex-complete}
You are told that $\vec{x}_p=\colvec{1\\1\\0}$ solves $\mat{A}\vec{x}=\vec{b}$ for
some matrix $\mat{A}$, and that $\nullsp(\mat{A})=\spn\{(1,-1,2)\}$. Write down the
complete solution set of $\mat{A}\vec{x}=\vec{b}$. Is the solution unique? Give
two distinct solutions explicitly.
\end{exercise}

\begin{exercise}[More unknowns than equations]\label{ch12:ex-wide}
Explain, in two or three sentences and citing the pivot count, why a homogeneous
system of $3$ equations in $5$ unknowns must always have a nonzero solution.
Construct one explicit $3\times5$ example and exhibit a nonzero nullspace vector.
\end{exercise}

\begin{exercise}[Subspace, true or false]\label{ch12:ex-tf}
For each set, decide whether it is a subspace and justify in one line.
(a) $\{\vec{x}\in\R^3 : \mat{A}\vec{x}=\vec{0}\}$ for a fixed $\mat{A}$.
(b) $\{\vec{x}\in\R^3 : \mat{A}\vec{x}=\vec{b}\}$ for a fixed $\vec{b}\neq\vec0$.
(c) $\col(\mat{A})$ for a $4\times2$ matrix, viewed as a subset of $\R^4$.
(d) The set of all $\vec{b}$ for which $\mat{A}\vec{x}=\vec{b}$ has \emph{no}
solution.
\end{exercise}

\begin{exercise}[Geometry of the collapse]\label{ch12:ex-collapse}
Let $\mat{A}=\begin{bmatrix}2 & 6\\1 & 3\end{bmatrix}$. Draw the unit square's
image under $\mat{A}$ and explain why it is a segment, not a parallelogram.
Identify the column space (the line the square collapses onto) and the nullspace
(the direction collapsed to the origin), and sketch both in the plane.
\end{exercise}

\begin{exercise}[Proof: nullspace under stacking]\label{ch12:ex-proof-stack}
Suppose $\mat{B}$ is the matrix obtained by stacking $\mat{A}$ on top of a copy of
itself, $\mat{B}=\begin{bmatrix}\mat{A}\\\mat{A}\end{bmatrix}$. Prove that
$\nullsp(\mat{B})=\nullsp(\mat{A})$. Then explain in words why this matches the
intuition that duplicating equations adds no new information.
\end{exercise}

\begin{exercise}[Challenge: nullspace of a product]\label{ch12:ex-product}
Let $\mat{A}$ be $m\times n$ and $\mat{B}$ be $n\times p$. Prove that
$\nullsp(\mat{B})\subseteq\nullsp(\mat{A}\mat{B})$. Give a $2\times2$ example
where the containment is \emph{strict} (the product kills a vector that $\mat{B}$
does not), and explain geometrically how $\mat{A}$ can collapse extra directions
that $\mat{B}$ left alive.
\end{exercise}
